\documentclass[10pt,a4paper]{article}
\usepackage{../../styles/neet_style}

\begin{document}

% Title Header
\begin{center}
  {\Huge \textbf{NEET UG 2027: Master Biology Chapter Notes}}\\[6pt]
  {\Large \textbf{Cell Biology, Cell Cycle, Animal Diversity \& Anatomy (Frog \& Cockroach)}}\\[4pt]
  {\small \textbf{Prepared, Compiled \& Authored by Satya Prakash} $\cdot$ Strictly Calibrated to Rationalized NCERT \& NMC Scope}
  \vspace{4pt}
  \hrule height 1.5pt
\end{center}

\pagestyle{fancy}
\fancyhf{}
\fancyhead[L]{\small \textbf{Cell Biology, Animal Diversity \& Morphology}}
\fancyhead[R]{\small Satya Prakash}
\fancyfoot[L]{\small \textbf{NEET UG 2027 Prep Series}}
\fancyfoot[C]{\thepage}
\fancyfoot[R]{\small \textbf{Zoology Chapter Notes}}

\section{Cell Structure, Function \& Biomolecules}

\begin{conceptbox}[Prokaryotic vs Eukaryotic Cell Architecture]
\textbf{1. Cell Theory Milestones:}
\begin{itemize}
  \item \textbf{M.J. Schleiden (1838)} observed plants; \textbf{Theodore Schwann (1839)} studied animal cells and noted presence of cell wall unique to plant cells.
  \item \textbf{Rudolf Virchow (1855)} modified cell theory: \textbf{\textit{Omnis cellula-e cellula}} (all cells arise from pre-existing cells).
\end{itemize}
\textbf{2. Prokaryotic Diagnostic Architecture:}
\begin{itemize}
  \item \textbf{Cell Envelope (3 Layers):} Outermost \textbf{Glycocalyx} (loose slime layer or rigid \textbf{capsule}) $\to$ Middle \textbf{Peptidoglycan cell wall} $\to$ Innermost semi-permeable \textbf{Plasma membrane}.
  \item \textbf{Mesosomes:} Invaginations of plasma membrane in the form of vesicles, tubules, and lamellae; increase surface area for enzymatic respiration, secretion, and replication.
  \item \textbf{Plasmids:} Small, circular, self-replicating extrachromosomal dsDNA; confer specific phenotypes like antibiotic resistance.
  \item \textbf{Inclusion Bodies:} Non-membrane bound storage granules free in cytoplasm (e.g., phosphate granules, cyanophycean granules, glycogen granules, gas vacuoles in blue-green and purple photosynthetic bacteria).
\end{itemize}
\textbf{3. Eukaryotic Organelles \& The Endomembrane System:}
The endomembrane system includes functionally coordinated organelles: \textbf{Endoplasmic Reticulum (ER), Golgi Apparatus, Lysosomes, and Vacuoles} (Mitochondria, Chloroplasts, and Peroxisomes are NOT part of it).
\begin{itemize}
  \item \textbf{Fluid Mosaic Model (Singer \& Nicolson, 1972):} Quasi-fluid lipid bilayer (phospholipids with hydrophilic heads outside, hydrophobic fatty acid tails inside) with integrated peripheral and integral proteins.
  \item \textbf{Rough ER (RER):} Studded with $80S$ ribosomes; active in \textbf{protein synthesis and secretion}.
  \item \textbf{Smooth ER (SER):} Major site for \textbf{lipid and steroid hormone synthesis}; detoxifies drugs in liver.
  \item \textbf{Golgi Complex (Camillo Golgi, 1898):} Flattened cisternae ($0.5-1.0\,\mu\mathrm{m}$); convex forming \textbf{cis-face} receives vesicles from ER; concave maturing \textbf{trans-face} releases modified vesicles; crucial site for \textbf{glycosylation of proteins (glycoproteins)} and \textbf{lipids (glycolipids)}.
  \item \textbf{Lysosomes:} Formed by Golgi packaging; contain hydrolytic enzymes (lipases, proteases, carbohydrases) active at \textbf{acidic pH} ($\approx 5.0$).
  \item \textbf{Ribosomes (George Palade, 1953):} Non-membrane bound ribonucleoprotein granules ($70S = 50S + 30S$ in prokaryotes/organelles; $80S = 60S + 40S$ in eukaryotic cytoplasm).
  \item \textbf{Cilia \& Flagella:} Possess $9+2$ axoneme microtubules; basal body derived from centrioles ($9+0$ cartwheel).
\end{itemize}
\end{conceptbox}

\begin{conceptbox}[Enzyme Catalysis, Kinetics \& Michaelis-Menten Dynamics]
Enzymes are biocatalysts (mostly proteins; ribozymes are RNA enzymes like $23S$ rRNA):
\begin{itemize}
  \item Lower the \textbf{activation energy} barrier without altering standard free energy change ($\Delta G$) or equilibrium constant ($K_{\mathrm{eq}}$).
  \item \textbf{Michaelis Constant ($K_m$):} Substrate concentration at which the reaction velocity reaches \textbf{half-maximal velocity} ($V_{\max}/2$).
  \begin{equation}
    K_m \propto \frac{1}{\text{Substrate Affinity of Enzyme}} \quad (\text{Lower } K_m \implies \text{Higher catalytic affinity})
  \end{equation}
  \item \textbf{Competitive Inhibition:} Inhibitor closely resembles substrate in molecular structure and competes for active site:
  \begin{equation}
    \textbf{Competitive Inhibitor Effect:} \quad \mathbf{K_m \text{ increases}} \qquad \mathbf{V_{\max} \text{ remains unchanged}}
  \end{equation}
  \textit{Canonical NCERT Example:} Inhibition of \textbf{succinate dehydrogenase by malonate} (resembles succinate).
  \item \textbf{Non-Competitive Inhibition:} Inhibitor binds to allosteric site: $\mathbf{V_{\max} \text{ decreases}}$, while $\mathbf{K_m \text{ remains unchanged}}$.
  \item \textbf{Enzyme Classes (IUBMB - 6 Classes):} (1) \textbf{Oxidoreductases}, (2) \textbf{Transferases}, (3) \textbf{Hydrolases}, (4) \textbf{Lyases} (cleave bonds without water), (5) \textbf{Isomerases}, (6) \textbf{Ligases} (join bonds with ATP hydrolysis).
  \item \textbf{Cofactors (Non-Protein Moieties):}
  \begin{itemize}
    \item \textbf{Prosthetic Groups:} Tightly and permanently bound organic compounds (e.g., \textbf{Haem} in peroxidase and catalase).
    \item \textbf{Coenzymes:} Transiently bound organic compounds derived from vitamins (e.g., \textbf{NAD and NADP} contain niacin).
    \item \textbf{Metal Ions:} Form coordination bonds (e.g., \textbf{$\mathrm{Zn}^{2+}$} in carboxypeptidase and carbonic anhydrase).
  \end{itemize}
\end{itemize}
\end{conceptbox}

\section{Cell Cycle \& Cell Division: Mitosis \& Meiosis}

\begin{formulabox}[Cell Cycle Progression \& Chromosome/DNA Mathematics]
Standard eukaryotic cell cycle (human cell divides once in $\approx 24$ hours; yeast divides in $90$ minutes):
\begin{align*}
  \textbf{Interphase ($>95\%$ of cycle): } & \mathbf{G_1} \ (\text{growth}) \longrightarrow \mathbf{S} \ (\text{DNA replication}) \longrightarrow \mathbf{G_2} \ (\text{spindle proteins})\\
  \textbf{M-Phase ($<5\%$ of cycle): } & \text{Prophase} \longrightarrow \text{Metaphase} \longrightarrow \text{Anaphase} \longrightarrow \text{Telophase} \longrightarrow \text{Cytokinesis}
\end{align*}
\textbf{Quantitative Dynamics of Chromosome Number ($2n$) and DNA Content ($C$):}
\begin{center}
\scriptsize
\begin{tabular}{|l|c|c|c|c|c|c|}
\hline
\textbf{Stage} & \textbf{$G_1$ Phase} & \textbf{S Phase} & \textbf{$G_2$ Phase} & \textbf{Metaphase} & \textbf{Anaphase} & \textbf{Daughter Cells} \\
\hline
Chromosome Count & $2n$ & $2n$ & $2n$ & $2n$ & $\mathbf{4n}$ (split centromeres) & $2n$ \\
DNA Content & $2C$ & $\mathbf{4C}$ (doubles) & $4C$ & $4C$ & $4C$ & $2C$ \\
\hline
\end{tabular}
\end{center}
\end{formulabox}

\begin{conceptbox}[Meiotic Substages: Prophase I Molecular Events]
Meiosis reduces chromosome number by half ($2n \to n$). Prophase I is prolonged and divided into 5 distinct substages:
\begin{enumerate}
  \item \textbf{Leptotene:} Chromosomes gradually become visible under light microscope; compaction of chromosomes continues.
  \item \textbf{Zygotene:} Homologous chromosomes pair together (\textbf{Synapsis}), forming a nucleoprotein complex called \textbf{Synaptonemal Complex}. The paired complex is termed a \textbf{Bivalent} or \textbf{Tetrad}.
  \item \textbf{Pachytene:} Non-sister chromatids of homologous chromosomes exchange genetic material (\textbf{Crossing Over}) at recombination nodules, catalyzed by the enzyme \textbf{Recombinase}.
  \item \textbf{Diplotene:} Dissolution of synaptonemal complex; homologous chromosomes separate except at crossover sites, forming X-shaped configurations called \textbf{Chiasmata}. (Can remain arrested for months/years in oocytes of vertebrates, e.g., Dictyotene).
  \item \textbf{Diakinesis:} \textbf{Terminalization of chiasmata}; chromosomes are fully condensed; nucleolus disappears and nuclear envelope breaks down.
\end{enumerate}
\textbf{Crucial Meiotic Checkpoints:}
\begin{itemize}
  \item \textbf{Anaphase I:} Homologous chromosomes separate to opposite poles; \textbf{sister chromatids remain attached at their centromeres} $\implies$ Reduction in chromosome number ($2n \to n$).
  \item \textbf{Anaphase II:} \textbf{Centromeres split}, allowing sister chromatids to move to opposite poles $\implies$ Resembles normal mitotic anaphase.
\end{itemize}
\end{conceptbox}

\section{Animal Kingdom: Non-Chordates \& Chordates Blueprint}

\begin{conceptbox}[Diagnostic Hallmarks of Major Animal Phyla]
\begin{center}
\scriptsize
\begin{tabular}{|p{2.3cm}|p{2.3cm}|p{2.6cm}|p{7.0cm}|}
\hline
\textbf{Phylum} & \textbf{Level of Org. / Symmetry} & \textbf{Coelom / Germ Layers} & \textbf{Diagnostic Features \& Canonical Examples} \\
\hline
\textbf{Porifera} (Sponges) & Cellular; Asymmetrical & Acoelomate; Diploblastic & \textbf{Choanocytes / Collar cells} line spongocoel; Water canal system; Spicules/spongin skeleton. \textit{Sycon, Spongilla} (freshwater), \textit{Euspongia} (bath sponge). \\
\hline
\textbf{Cnidaria} (Coelenterata) & Tissue; Radial symmetry & Acoelomate; Diploblastic & \textbf{Cnidoblasts / Nematocysts} on tentacles for defense/anchorage; Polyp and Medusa; \textbf{Metagenesis} in \textit{Obelia}. \textit{Hydra, Physalia} (Portuguese man-of-war), \textit{Adamsia} (sea anemone). \\
\hline
\textbf{Ctenophora} (Comb jellies) & Tissue; Radial symmetry & Acoelomate; Diploblastic & Eight external rows of ciliated \textbf{comb plates} for locomotion; \textbf{Bioluminescence}; Strictly marine. \textit{Pleurobrachia, Ctenoplana}. \\
\hline
\textbf{Platyhelminthes} (Flatworms) & Organ / Organ-system; Bilateral & \textbf{Acoelomate}; Triploblastic & Dorso-ventrally flattened body; Hooks and suckers in parasitic forms; \textbf{Flame cells (Protonephridia)} for osmoregulation. \textit{Taenia} (tapeworm), \textit{Fasciola} (liver fluke). \\
\hline
\textbf{Aschelminthes} (Roundworms) & Organ-system; Bilateral & \textbf{Pseudocoelomate}; Triploblastic & Circular cross-section; Muscular pharynx; Sexual dimorphism (female longer). \textit{Ascaris, Wuchereria} (filarial worm), \textit{Ancylostoma} (hookworm). \\
\hline
\textbf{Annelida} (Segmented worms) & Organ-system; Bilateral & \textbf{Coelomate}; Triploblastic & \textbf{Metameric segmentation}; Longitudinal/circular muscles; \textbf{Nephridia} for excretion; Closed circulation. \textit{Nereis} (parapodia), \textit{Pheretima} (earthworm), \textit{Hirudinaria} (leech). \\
\hline
\textbf{Arthropoda} (Jointed legs) & Organ-system; Bilateral & Coelomate; Triploblastic & \textbf{Largest phylum} ($>2/3$ species); Chitinous exoskeleton; \textbf{Malpighian tubules} for excretion; Statocysts; Open circulation. \textit{Apis, Bombyx, Laccifer, Anopheles, Locusta, Limulus} (living fossil). \\
\hline
\textbf{Mollusca} (Soft-bodied) & Organ-system; Bilateral & Coelomate; Triploblastic & \textbf{Second largest phylum}; Calcareous shell; Unsegmented with head, muscular foot, and visceral hump; Mantle cavity with ctenidia; Rasping organ \textbf{Radula}. \textit{Pila, Pinctada, Sepia, Octopus}. \\
\hline
\textbf{Echinodermata} (Spiny-skinned) & Organ-system; Radial (adult), Bilateral (larva) & Coelomate; Triploblastic & Endoskeleton of calcareous ossicles; \textbf{Water vascular system} (ambulacral system) for locomotion, food capture, and respiration; Excretory system absent. \textit{Asterias} (starfish), \textit{Echinus, Antedon}. \\
\hline
\textbf{Hemichordata} & Organ-system; Bilateral & Coelomate; Triploblastic & Worm-like marine; Proboscis, collar, and trunk; \textbf{Stomochord} in collar; Open circulation; Proboscis gland for excretion. \textit{Balanoglossus, Saccoglossus}. \\
\hline
\textbf{Chordata} & Organ-system; Bilateral & Coelomate; Triploblastic & Fundamental hallmarks: (1) \textbf{Notochord}, (2) \textbf{Dorsal hollow nerve cord}, (3) \textbf{Paired pharyngeal gill slits}, (4) \textbf{Post-anal tail}. \\
\hline
\end{tabular}
\end{center}
\end{conceptbox}

\section{Structural Organisation in Animals: Frog \& Cockroach Anatomy}

\begin{conceptbox}[Cockroach Anatomy (\textit{Periplaneta americana}) Prescribed Scope]
\begin{itemize}
  \item \textbf{Morphology:} Body divided into Head, Thorax, and Abdomen; covered by brown chitinous exoskeleton with hardened plates called \textbf{Sclerites} (dorsal \textbf{Tergites}, ventral \textbf{Sternites}, joined by flexible \textbf{arthrodial membrane}).
  \item \textbf{Head:} Triangular, hypognathous; formed by fusion of \textbf{6 embryonic segments}; bears a pair of compound eyes (each composed of $\approx \mathbf{2000\text{ hexagonal ommatidia}}$ providing mosaic vision with high sensitivity but low resolution).
  \item \textbf{Mouthparts (Biting and Chewing Type):} Labrum (upper lip), paired Mandibles (with grinding and incising teeth), paired Maxillae, Labium (lower lip), and median flexible lobe \textbf{Hypopharynx} (tongue).
  \item \textbf{Thorax:} Prothorax, Mesothorax, Metathorax. Forewings (\textbf{Tegmina} from mesothorax) are opaque, dark, leathery, and protective at rest; Hindwings (from metathorax) are transparent, membranous, and used in \textbf{flight}.
  \item \textbf{Abdomen:} Consists of \textbf{10 segments} in both sexes.
  \begin{itemize}
    \item \textbf{Sexual Dimorphism:} Males bear a pair of unjointed \textbf{Anal Styles} on the $9^{\mathrm{th}}$ sternite (absent in females). Both sexes bear a pair of jointed \textbf{Anal Cerci} on the $10^{\mathrm{th}}$ segment.
  \end{itemize}
  \item \textbf{Digestive System:} Alimentary canal: Mouth $\to$ Pharynx $\to$ Oesophagus $\to$ \textbf{Crop} (storage) $\to$ \textbf{Gizzard / Proventriculus} (has outer circular muscles and 6 thick chitinous teeth for grinding food).
  \begin{itemize}
    \item \textbf{Hepatic / Gastric Caeca:} $6-8$ blind tubules at the junction of foregut and midgut; secrete digestive juice.
    \item \textbf{Malpighian Tubules:} $100-150$ thin, yellow, filamentous tubules at the junction of midgut and hindgut; extract uric acid from hemolymph (\textbf{Uricotelic}).
  \end{itemize}
  \item \textbf{Circulatory System:} Open type; heart consists of \textbf{13 funnel-shaped chambers} with ostia on either side.
  \item \textbf{Nervous System:} Consists of a series of fused ganglia: \textbf{3 in the thorax and 6 in the abdomen}. If head is cut off, cockroach survives for up to a week because the head holds only a bit of the nervous system!
\end{itemize}
\end{conceptbox}

\begin{conceptbox}[Frog Anatomy (\textit{Rana tigrina}) Prescribed Scope]
\begin{itemize}
  \item \textbf{Morphology:} Cold-blooded (\textbf{Poikilotherm}); amphibious; smooth, moist, slimy skin with mucus glands; camouflages (\textbf{Mimicry}). Body divided into \textbf{Head and Trunk} (neck and tail absent). Eyes are protected by a \textbf{Nictitating membrane}. Tympanum receives sound vibrations.
  \begin{itemize}
    \item \textbf{Sexual Dimorphism:} Male frogs have sound-producing \textbf{Vocal Sacs} and a \textbf{Copulatory pad} on the first digit of the forelimb (both absent in female).
  \end{itemize}
  \item \textbf{Respiration:} Cutaneous (skin in water and hibernation/aestivation), Buccopharyngeal cavity, and Pulmonary (lungs on land).
  \item \textbf{Circulatory System:} Closed, 3-chambered heart (\textbf{2 Atria and 1 Ventricle}).
  \begin{itemize}
    \item Triangular \textbf{Sinus Venosus} joins the right atrium; receives blood from vena cava.
    \item Ventricle opens ventrally into \textbf{Conus Arteriosus}.
    \item \textbf{RBCs:} Nucleated, oval, and contain hemoglobin (unlike non-nucleated human RBCs).
    \item Well-developed \textbf{Hepatic Portal} and \textbf{Renal Portal} systems.
  \end{itemize}
  \item \textbf{Excretory System:} Paired compact, dark red \textbf{Mesonephric kidneys}; \textbf{Ureotelic} (excretes urea). Functional unit is uriniferous tubule. In males, ureters act as \textbf{Urogenital ducts} opening into the \textbf{Cloaca} (common chamber for digestive, excretory, and reproductive products).
  \item \textbf{Male Reproductive System:} Paired yellowish ovoid testes adhered to upper kidneys by double fold of peritoneum called \textbf{Mesorchium}. $10-12$ \textbf{Vasa efferentia} enter kidneys to run into \textbf{Bidder's Canal}, opening into urogenital duct.
  \item \textbf{Development:} External fertilization in water; indirect development via swimming aquatic larva called \textbf{Tadpole} (undergoes metamorphosis governed by thyroxine).
\end{itemize}
\end{conceptbox}

\end{document}
